\section*{第 \ref{ch:b2:genome-diversification} 章——突变与基因组多样化}

\begin{solution}{exo:b2:genome-diversification:1}
GAG：沉默\omterm{def:b2:genome-diversification:mutation}{突变}（仍为谷氨酸），转换（A$\to$G）。GTA：错义
（缬氨酸），颠换（A$\to$T）。TAA：无义（终止），
颠换（G$\to$T）。插入一个碱基：\omterm{def:b2:genome-diversification:mutation}{移码}，
之后的每个密码子都会被改写。
\end{solution}

\begin{solution}{exo:b2:genome-diversification:2}
$U = 5\times 10^{-10}\times 4.6\times 10^{6} = 2.3\times 10^{-3}$：约每
430 次分裂产生一个\omterm{def:b2:genome-diversification:mutation}{突变}。在 $10^{9}$ 个细胞中：
$2.3\times 10^{6}$ 个新\omterm{def:b2:genome-diversification:mutation}{突变}每代。
\end{solution}

\begin{solution}{exo:b2:genome-diversification:3}
转化：从裂解细胞摄取游离 DNA，需要有感受态的受体；
转移染色体片段（肺炎链球菌中为荚膜基因）。接合：细胞间
通过菌毛接触，需要供体中有接合性\omterm{def:b2:unicellular-diversity:prokaryote}{质粒}；
转移\omterm{def:b2:unicellular-diversity:prokaryote}{质粒}（耐药基因），在 Hfr 菌株中还可转移染色体基因。
\omterm{def:b2:genome-diversification:hgt}{转导}：噬菌体包装了宿主 DNA；转移随机片段
（普遍性\omterm{def:b2:genome-diversification:hgt}{转导}），或紧邻原噬菌体整合位点的基因
（专一性\omterm{def:b2:genome-diversification:hgt}{转导}）。
\end{solution}

\begin{solution}{exo:b2:genome-diversification:4}
\omterm{def:b2:genome-diversification:polyploidy}{同源多倍体}：来自同一物种的额外染色体组，由二倍体配子产生
（四倍体马铃薯，许多同源四倍体三叶草）。\omterm{def:b2:genome-diversification:polyploidy}{异源多倍体}：两个物种的
染色体组经杂交合并后再加倍（面包小麦、棉花、烟草）。
二倍体杂种 AB 中没有两条相同的染色体，因此
减数第一次分裂时无法配对；染色体随机分离，
配子不平衡且不可育。
\end{solution}

\begin{solution}{exo:b2:genome-diversification:5}
$2\times 3.2\times 10^{9}\times 1.2\times 10^{-8} = 77$ 个新\omterm{def:b2:genome-diversification:mutation}{突变}。
编码序列中：$77\times 0.015 = 1.2$；改变氨基酸：
$1.2\times 0.7 = 0.8$——大多数
儿童携带约一个新的氨基酸替换。
\end{solution}

\begin{solution}{exo:b2:genome-diversification:6}
\omterm{def:b2:genome-diversification:mutation}{突变}事件数：$m(N - N_0) = 2\times 10^{-8}\times 10^{9} = 20$。\omterm{def:b2:genome-diversification:mutation}{突变}细胞数：
$mN\ln(N/N_0) = 20\times \ln 10^{6} = 20\times 13.8 = 276$。
$p_0 = e^{-20} \approx 2\times 10^{-9}$：每个培养物中都有\omterm{def:b2:genome-diversification:mutation}{突变}体，因此
统计不含\omterm{def:b2:genome-diversification:mutation}{突变}体的培养物则什么也测不出。需要 $m(N -
N_0) \approx 1$，即 $N \approx 5\times 10^{7}$，此时 $p_0 \approx
0.37$，数十个培养物即可将 $m$ 测定到两倍以内。
\end{solution}

\begin{solution}{exo:b2:genome-diversification:7}
DNA 中本不应有尿嘧啶，因此尿嘧啶-DNA 糖基化酶
可以无歧义地去除它发现的每一个尿嘧啶。胸腺嘧啶是正常碱基：
5-甲基胞嘧啶脱氨后出现 G$\cdot$T 错配，
但修复机制无法判断哪条链是错误的，
有一半的情况下会“纠正”G，从而固定 C$\to$T 转换。
因此甲基化的 CpG 位点\omterm{thm:b2:genome-diversification:rate}{突变速率}约为其他位点的十倍，
经过长期进化，它们已转变为 TpG 和 CpA，
哺乳动物基因组中所含 CpG
仅约为碱基组成所预测数量的五分之一。
\end{solution}

\begin{solution}{exo:b2:genome-diversification:8}
\qty{46}{kb} 每分钟。位置：368、782、1150 和
\qty{1840}{kb}；间距 414、368 和 \qty{690}{kb}。
\end{solution}

\begin{solution}{exo:b2:genome-diversification:9}
一条染色体的第二个拷贝与另一条染色体的第一个拷贝配对；二者之间的
交换产生一条染色单体，携带拷贝~1、拷贝~2 和拷贝~2$'$（三个基因），另一条
染色单体仅携带拷贝~1$'$。若某人从每个亲本各获得一条只含一个基因的染色体，
则其只有两个 $\alpha$ 基因，而非四个（$-\alpha/-\alpha$）：$\alpha$
链以一半速率合成，红细胞小而
色淡，贫血程度轻微或无贫血——即
$\alpha$-地中海贫血性状。
\end{solution}

\begin{solution}{exo:b2:genome-diversification:10}
配子：$n = 21$（A、B、D 各一个染色体组）。杂种：AB，14 条染色体，
ABD，21 条。面包小麦 $\times$ 二粒小麦：AABBD，35 条染色体——
A 与 B 染色体组可以配对，D 染色体组没有配对伙伴，
植株基本不育。每次加倍使每条染色体都有一条完全相同的同源染色体，
因此形成二价体，配子平衡；
与亲本回交则产生一个染色体组不配对的植株，
其配子不平衡，因此该六倍体与其祖先生殖隔离——
一步即成一个新物种。
\end{solution}

\begin{solution}{exo:b2:genome-diversification:11}
逻辑斯谛增长：$R(t) = N/\bigl(1 + (N - 1)e^{-\gamma Nt}\bigr)$。
当 $(N - 1)e^{-\gamma Nt} = 1$ 时种群达到一半：$t = \ln(N -
1)/\gamma N = 20.7/0.5 = \qty{41}{h}$。在无抗生素时，
不含\omterm{def:b2:unicellular-diversity:prokaryote}{质粒}的细胞生长快 \qty{5}{\%}，因此携带者与非携带者之比
以 $e^{-st}$ 衰减，其中 $s = 0.05\ln 2 =
\qty{0.035}{h^{-1}}$（\omterm{def:b2:unicellular-diversity:division}{世代时间} \qty{1}{h}）：从 \qty{99}{\%} 降至
\qty{1}{\%} 需要 $\ln(99^2)/0.035 = \qty{260}{h}$，
约十一天——而持续的接合转移可使\omterm{def:b2:unicellular-diversity:prokaryote}{质粒}
以较低频率长期保留，为下一疗程的药物
做好准备。
\end{solution}

\begin{solution}{exo:b2:genome-diversification:12}
\omterm{thm:b2:genome-diversification:luria}{波动试验}表明，耐药\omterm{def:b2:genome-diversification:mutation}{突变}在选择因素出现之前、与之无关地、
随机的时间发生，因此\omterm{def:b2:genome-diversification:mutation}{突变}对何种变异有用是“盲目”的；
影印平板实验在完全不接触被选择细胞的情况下得到了相同的结论。
进化并非随机，因为自然选择按照变异的后果
对这些盲目产生的变异进行筛选，且筛选是持续、累积的。
唯一的随机性本身受到塑造之处，
在于\omterm{thm:b2:genome-diversification:rate}{突变速率}：错误过多或过少的谱系都会被淘汰，
因此\omterm{thm:b2:genome-diversification:rate}{突变速率}是一种经过适应的折中——
但每一次\omterm{def:b2:genome-diversification:mutation}{突变}造成什么后果依然不可预测，也不经挑选。
\end{solution}

\begin{solution}{pb:b2:genome-diversification:1}
\textbf{1.} 和为 119，均值 $5.95$；平方均值 $11489/20 = 574$；
方差 $574 - 35 = 540$。
\textbf{2.} 和为 324，均值 $16.2$；平方均值 $5462/20 = 273$；
方差 $273 - 262 = 11$。
\textbf{3.} 同一培养物的样本服从泊松分布（方差接近均值）：
它们是从同一固定\omterm{def:b2:genome-diversification:mutation}{突变}体比例中随机抽取的。
独立培养物的方差为均值的九十倍：
其\omterm{def:b2:genome-diversification:mutation}{突变}体在生长过程中于不同的随机时刻出现，
\omterm{def:b2:genome-diversification:mutation}{突变}出现越早，其克隆越大。
\textbf{4.} 大奖：\omterm{def:b2:genome-diversification:mutation}{突变}发生在最初几次分裂之一，
其克隆随培养物一同生长至一百个细胞。
\textbf{5.} 二十个中十四个为空：$p_0 = 0.70$；$m(N - N_0) = -\ln 0.70
= 0.36$；$m = 0.36/2\times 10^{8} = \qty{1.8e-9}{}$ 每次分裂。
\textbf{6.} $1.8\times 10^{-9}\times 2\times 10^{8}\times \ln(2\times
10^{6}) = 0.36\times 14.5 = 5.2$ 个\omterm{def:b2:genome-diversification:mutation}{突变}细胞；观测值 $5.95$——
与二十个培养物的随机波动范围内吻合。
\textbf{7.} 没有样本为空（所有样本都共享亲代培养物的
\omterm{def:b2:genome-diversification:mutation}{突变}体），因此 $p_0 = 0$ 无法给出估计值；这些样本测量的是
同一培养物的\omterm{def:b2:genome-diversification:mutation}{突变}体比例，该比例取决于其\omterm{def:b2:genome-diversification:mutation}{突变}
恰好发生的时刻。
\textbf{8.} $u = m/20 = \qty{9e-11}{}$ 每碱基对每次复制。
\textbf{9.} $U = 9\times 10^{-11}\times 4.6\times 10^{6} =
\qty{4.1e-4}{}$。
\textbf{10.} $10^{12}$ 个细胞 $\times$ $4.1\times 10^{-4}$ $= 4.1\times
10^{8}$ 个新\omterm{def:b2:genome-diversification:mutation}{突变}。
\textbf{11.} $3\times 4.6\times 10^{6} = 1.4\times 10^{7}$ 种可能的
改变；每种改变出现 $4.1\times 10^{8}/1.4\times 10^{7} \approx 30$
次。
\textbf{12.} 每一种单碱基改变，包括每一种
赋予细胞对某种一次\omterm{def:b2:genome-diversification:mutation}{突变}即可克服的药物耐药性的改变，均在加入药物前便已存在于
数十个细胞中；药物只是
对它们进行筛选。
\textbf{13.} 同一细胞内两次独立\omterm{def:b2:genome-diversification:mutation}{突变}的发生率为
$m^2 \approx (2\times 10^{-9})^2 = 4\times 10^{-18}$ 每次分裂：在
$10^{12}$ 个细胞中为 $4\times 10^{-6}$ 每代——双重
耐药细胞几乎从不预先存在，这正是
结核病需同时用多种药物治疗的原因。
\textbf{14.} 令 $a = N - 1$，$E = e^{-\gamma Nt}$：$R = N/(1 +
aE)$，$\mathrm{d}R/\mathrm{d}t = \gamma N\cdot NaE/(1 + aE)^2$；且
$\gamma R(N - R) = \gamma\,\frac{N}{1 + aE}\cdot\frac{NaE}{1 + aE}$，
两者相同；$R(0) = N/(1 + N - 1) = 1$。
\textbf{15.} $(N - 1)E = 1$：$t = \ln(N - 1)/\gamma N = 20.7/0.5 =
\qty{41}{h}$。
\textbf{16.} $1 + aE = 1/0.99$，$aE = 0.0101$：$t = \ln(10^{9}/0.0101)
/0.5 = (20.7 + 4.6)/0.5 = \qty{51}{h}$。
\textbf{17.} $\gamma (N/2)^2 = \gamma N\cdot N/4 = 0.5\times 2.5\times
10^{8} = 1.25\times 10^{8}$ 次转移每小时。
\textbf{18.} 增长速率分别为 $\ln 2 = \qty{0.693}{h^{-1}}$ 和
\qty{0.762}{h^{-1}}；\omterm{def:b2:genome-diversification:mutation}{突变}体与敏感型之比以 $e^{st}$
增长，其中 $s = \qty{0.069}{h^{-1}}$；从 $1/N$ 增至 1 所需时间为 $\ln N/s =
20.7/0.069 = \qty{300}{h}$，即十二天。
\textbf{19.} 转移正比于携带者与非携带者的乘积，
且以接触速率而非分裂速率进行；
扩散受限于\omterm{def:b2:genome-diversification:mutation}{突变}体的生长优势，而该优势很小。
\omterm{def:b2:unicellular-diversity:prokaryote}{质粒}在两天内即可遍及种群，而有利\omterm{def:b2:genome-diversification:mutation}{突变}体
需两周。
\textbf{20.} $U = 4.1\times 10^{-2}$；有害\omterm{def:b2:genome-diversification:mutation}{突变}：$0.041\times 0.88\times
0.3 = 0.011$ 每代。
\textbf{21.} 均值 $2.2$：$e^{-2.2} = 0.11$，即每九个后代中有一个
未受损伤。
\textbf{22.} 野生型：$1.1\times 10^{-4}$ 每代，
200 代内为 $0.022$：$e^{-0.022} = 0.98$。
\textbf{23.} 均值为 $520$ 个\omterm{def:b2:genome-diversification:mutation}{突变}细胞；$p_0 = e^{-36} \approx
10^{-16}$：不存在空培养物。
\textbf{24.} 在强烈且不断变化的选择下（如医院轮换使用抗生素），
新变异的百倍供应超过了
负荷，\omterm{def:b2:genome-diversification:mutation}{突变}增强株借其所产生的耐药性搭便车。在
稳定环境中无利可图，且
每百代一次有害\omterm{def:b2:genome-diversification:mutation}{突变}即会受损，因此\omterm{def:b2:genome-diversification:mutation}{突变}增强谱系
被清除。
\textbf{25.} 对耐药性，$m = \qty{1.8e-9}{}$ 每次分裂，$u =
\qty{9e-11}{}$ 每碱基对每次复制，$U = \qty{4.1e-4}{}$
每个基因组每次复制的\omterm{def:b2:genome-diversification:mutation}{突变}数；\omterm{def:b2:unicellular-diversity:prokaryote}{质粒}在
\qty{41}{h} 内到达肠道一半，并在 \qty{51}{h} 内达到其 \qty{99}{\%}。
\end{solution}
